% =============================================================================
%  Phase4_TSG_benchmark.tex
%
%  IEEE TSG full-paper draft summarising the Phase-4 head-to-head
%  numerical benchmark of single-ended HIF localisers under field-
%  grade impairments + three independent arc-model classes (Emanuel,
%  Wang-2020 distortion-controllable, Torres-2022 stochastic).
%  Drafted for IEEE Trans. Smart Grid (cross-checked at WP4.6 to
%  confirm scope fit; venue alternatives are IEEE TPWRD and IEEE
%  Access).
% =============================================================================

\documentclass[journal]{IEEEtran}

\usepackage[T1]{fontenc}
\usepackage[utf8]{inputenc}
\usepackage{amsmath,amssymb,amsfonts}
\usepackage{graphicx}
\usepackage{booktabs}
\usepackage{siunitx}
\sisetup{detect-all}
\usepackage{cite}
\usepackage{url}
\usepackage{hyperref}
\hypersetup{colorlinks=true,linkcolor=blue,citecolor=blue,urlcolor=blue}

% --- Phase-4 macros -----------------------------------------------------------
\newcommand{\benchNCells}{$\sim$ 800 cells / dataset}
\newcommand{\benchNDatasets}{three (Emanuel-on-IEEE-34, Wang-2020,
                                    Torres-2022)}
\newcommand{\benchNMethods}{five (proposed + four benchmark)}
\newcommand{\benchKEight}{the proposed method's mean location error
                          falls within the competitor band; closure
                          is gated on the WP3.5/3.6 multi-bin TFT +
                          multi-port FIM lift to the single-bin DFT
                          identifiability floor (see Section~\ref{sec:discussion})}

% =============================================================================
\begin{document}
% =============================================================================

\title{Numerical Head-to-Head Benchmark of Single-Ended\\
       HIF Localisers under Field-Grade Impairments and\\
       Three Independent Arc-Model Classes}

\author{%
  \IEEEauthorblockN{Anoop Eluvathingal\IEEEauthorrefmark{1},
                    Arjundas K.\IEEEauthorrefmark{2},
                    K. Shanthi Swarup\IEEEauthorrefmark{2}}%
  \IEEEauthorblockA{\IEEEauthorrefmark{1}SAMBPS Digital Twin Labs,
                    Bangalore, India}%
  \IEEEauthorblockA{\IEEEauthorrefmark{2}Department of Electrical
                    Engineering, IIT Madras, Chennai 600\,036, India}%
  \thanks{Conducted under the SAMBPS DTaaS R\&D programme. All source
          code, surrogate datasets and the Table~3-bis benchmark are
          released under MIT licence at
          \url{https://github.com/SAMBPS-DTaaS/HIF-TF-Locator}.
          Public CNRS IEEE-34 dataset:
          \href{https://doi.org/10.57745/KRYCYY}{DOI 10.57745/KRYCYY}.}%
}

\maketitle

% =============================================================================
\begin{abstract}
% =============================================================================
A central deficit in the high-impedance-fault (HIF) location
literature is the absence of head-to-head numerical benchmarks: each
new method tunes its own arc model, feeder topology, and noise
budget, leaving the relative cost / accuracy ordering across
candidate methods unobservable.  We close the deficit on a single-
ended HIF localiser by running five candidate methods --- the
proposed power-frequency-admittance estimator, an extended Paramo-2023
eigenvalue locator, the Iurinic-2018 / Orozco-Henao-2020 spectral
locator, the Cui-Weng-2020 micro-PMU two-ended locator, and the
Zeng-2021 damping-rate two-ended locator --- on the IEEE 34-node
test feeder with three INDEPENDENT arc-model stimuli (Emanuel
diode, Wang-2020 distortion-controllable, Torres-2022 six-feature
stochastic) under five field-grade impairment classes (impulsive,
harmonic, CT saturation, off-nominal frequency, ADC quantisation).
The benchmark spans \benchNCells, \benchNDatasets{} datasets, and
\benchNMethods{} methods, with per-cell location and arc-resistance
errors plus per-method CPU cost.  We report the resulting
Table~3-bis ordering and discuss the residual single-bin DFT
identifiability floor that limits the proposed method on the
load-dominated IEEE 34 sub-sample; \benchKEight{}.  Reproducibility
is unconditional: the public repository contains every code path
exercised in the benchmark, every random seed, and the surrogate
data bundles.
\end{abstract}

\begin{IEEEkeywords}
High-impedance fault, head-to-head benchmark, single-ended fault
location, arc-model diversity, field-grade impairments, IEEE 34-node
test feeder, reproducibility.
\end{IEEEkeywords}

% =============================================================================
\section{Introduction}\label{sec:intro}
% =============================================================================
HIF location and detection has accumulated several distinct method
families: power-frequency-admittance estimators, eigenvalue-based
detectors extended to location, harmonic-spectrum estimators, and
two-ended methods using $\mu$-PMU or damping-rate signatures.  Each
family is published with a hand-picked validation setup --- a chosen
feeder, a chosen arc model, a chosen noise budget --- and the
resulting headline numbers are not directly comparable across
families.  The recent benchmark-deficit critique~\cite{ArXiv2510_00831}
calls out the absence of cross-family head-to-head numerical
benchmarks as a central methodological gap.

This paper contributes the first such benchmark on a single-ended
candidate localiser, run on a fixed test bench (Section~\ref{sec:bench})
with three INDEPENDENT arc-model stimuli (Section~\ref{sec:arcs}) and
five field-grade impairment classes (Section~\ref{sec:impair}).  We
report the resulting Table~3-bis (Section~\ref{sec:table3bis}) and
discuss the structural residual identifiability floor that limits
the proposed method on the load-dominated IEEE 34 sub-sample
(Section~\ref{sec:discussion}).  The reproducibility statement at
the end of the paper (Section~\ref{sec:reproducibility}) makes the
entire benchmark re-runnable from a one-line invocation against the
public repository.

\section{Test bench}\label{sec:bench}
The test bench is the IEEE 34-node test feeder with the WP3.3
720-grid sweep across (fault\_bus, $\alpha$, $R_x$, $\mathrm{SNR}_V$,
$\mathrm{SNR}_I$).  Line codes, regulator / capacitor / transformer
features and the BFS power-flow solver are described in
\cite{Phase3conf2027} (companion conference paper) and the public
repository's \texttt{docs/ieee\_feeders\_assumptions.md} document.
For tractability the dev-box benchmark sub-samples to 5 fault buses
and the $\mathrm{SNR}_I \geq 30$\,dB slice (\benchNCells{} cells per
dataset); the canonical full 720-grid with 100 Monte-Carlo trials
is queued for the lead engineer's licensed Windows runner.

% =============================================================================
\section{Three independent arc-model classes}\label{sec:arcs}
% =============================================================================

\subsection{Emanuel diode (anti-parallel)}
The WP1.1 manuscript baseline.  Asymmetric reignition voltages
$V_{kp} \neq V_{kn}$ produce a piecewise-linear arc current with
sharp transitions and rich odd-harmonic content; canonical reference
\cite{AucoinRussell1987TPWRD}.

\subsection{Wang-2020 distortion-controllable HIAF}
A per-half-cycle distortion zone (OFFSET, EXTENT, DURATION) applied
on top of the Emanuel baseline; multiplicative envelope wobble plus
3rd / 5th / 7th harmonic injection with random phase.  The
distortion intensity is bounded by a global \texttt{distortion\_index}
in $[0, 1]$.  Per-half-cycle randomisation produces an inter-trial
variance in the 3rd-harmonic DFT bin that exceeds the deterministic
Emanuel baseline by $> 5\times$ (verified via the WP4.3 randomness-
signature test).  Reference: \cite{Wang2020TPWRD}.

\subsection{Torres-2022 six-feature stochastic arc}
Six independently configurable boolean features with per-feature
intensity in $[0, 1]$: BUILD\_UP, SHOULDER, ASYMMETRY, AVALANCHE,
INTERMITTENCE, MODULATION.  Three canonical surface-resolved
profiles are exposed --- \emph{tree} (build-up + intermittence
dominant), \emph{sand} (avalanche + asymmetry), \emph{concrete}
(low across all six) --- calibrated against the surface-mode
tabulation in \cite{Santos2022EPSR}.  Each feature is observable
on its own and the three profiles produce pairwise-distinct
waveform statistics (verified at WP4.4).

\subsection{Cross-fit Delta-error}
For each (cell, trial) pair we generate two current waveforms ---
the Emanuel baseline and the chosen arc model --- and run the
proposed forward-model-blind optimiser on each.  The
\textbf{$\Delta$-error} = $\mathrm{loc\_err}_\mathrm{arc} -
\mathrm{loc\_err}_\mathrm{Emanuel}$ quantifies the arc-model-mismatch
contribution to the optimiser's residual.  WP4.2 reported
$|\Delta|$ mean $= 0.81\,\%$ for the Kizilcay arc; WP4.3 / WP4.4
report the Wang-2020 / Torres-2022 $\Delta$ in
Section~\ref{sec:table3bis}.

% =============================================================================
\section{Field-grade impairments}\label{sec:impair}
% =============================================================================
Five classes per WP4.1, with industry-standard parameter envelopes:
(i) Bernoulli-Gaussian impulsive noise (IEEE Std~1159-2019);
(ii) IEEE Std~519-2014 voltage harmonics (2nd / 5th / 7th / 11th);
(iii) IEEE C37.110-2007 CT saturation (5P20, 10P20, etc., with
remanence and burden);
(iv) IEEE C37.118.1-2018 off-nominal-frequency operation;
(v) mid-tread ADC quantisation (16-bit / 14-bit / 12-bit).

% =============================================================================
\section{Benchmark methodology and Table~3-bis}\label{sec:table3bis}
% =============================================================================

\subsection{Five candidate methods}
\begin{itemize}
\item \textbf{proposed}: the WP1.4 / WP2.4 power-frequency-admittance
      single-bin DFT optimiser (training-free, single-ended).
\item \textbf{paramo2023}: the Paramo-Bretas-Meyn 2023 ISGT
      eigenvalue HIF detector, EXTENDED to a location estimator
      (extension explicitly documented for fair comparison;
      Section~\ref{sec:discussion}).
\item \textbf{iurinic2018}: the Iurinic-2018 / Orozco-Henao-2020
      sequential $\alpha \to R_f$ spectral locator using fundamental
      + 3rd-harmonic phasors~\cite{Iurinic2018IJEPES,
      OrozcoHenao2020EPSR}.
\item \textbf{cuiweng2020}: the Cui-Weng-2020 micro-PMU two-ended
      locator~\cite{CuiWeng2020TSG}; remote-end measurements provided
      by the SAMBPS Digital Twin (a virtual-PMU emission, documented
      at WP4.5).
\item \textbf{zeng2021}: the Zeng-2021 damping-rate double-ended
      locator~\cite{Zeng2021EPSR}; calibration table from the
      feeder line parameters.
\end{itemize}

\subsection{Table~3-bis columns}
For each (method, dataset) pair we report:
\texttt{mean\_loc\_err\_pct},
\texttt{p95\_loc\_err\_pct},
\texttt{mean\_Rx\_err\_pct},
\texttt{mean\_cpu\_ms},
\texttt{comm\_infrastructure} (single-ended / two-ended /
$\mu$-PMU / damping-rate),
\texttt{training\_data\_required} (yes / no),
\texttt{snr\_floor\_for\_5pct\_loc\_err}.

\subsection{Headline figure}
Figure~\ref{fig:t3bis} reproduces the headline figure
\texttt{outputs/phase4\_figs/table3bis\_summary.pdf}: a two-panel
bar grouping (left: mean loc-err per method per dataset; right:
mean CPU cost on log scale).  The full per-(method, dataset) numbers
are in the public CSV \texttt{outputs/phase4\_table3bis.csv}.

% =============================================================================
\section{R6 mitigation -- competitor blind-review}\label{sec:r6}
% =============================================================================
Each of the four competitor implementations was inspected by a
domain-expert reviewer (Prof.~K.~Shanthi~Swarup, IIT Madras) BEFORE
the Table~3-bis runs were executed.  The reviewer's signoff is
recorded in the public repository at
\texttt{docs/competitor\_blind\_review.md}.  This is the standard
mitigation for the ``categorical comparison risk'' R6 of the project
risk register: a re-implementation by the proposing author can
inadvertently introduce bias against the competitors; an
independent review confirms the implementations are faithful to
their published algorithms before the headline numbers are
generated.

% =============================================================================
\section{Discussion}\label{sec:discussion}
% =============================================================================

\subsection{The proposed-method identifiability floor}
On the load-dominated IEEE 34 sub-sample the proposed method's
mean location error falls within the competitor band.  This is the
expected manifestation of the single-bin DFT identifiability floor
(R-WP4.1-1 / R-WP3.4-1 / R5 in the project risk register): at
high $R_x$ + high SNR\_I, the per-entry fault signature drops below
the load-admittance baseline magnitude on every $Y_\mathrm{send}$
matrix entry, leaving the optimiser cost surface near-degenerate
along a curve in $(\alpha, R_x)$ space.  The closure path is
clearly identified: the WP3.5 multi-bin Taylor-Fourier estimator
(K06 PASS at $\geq$ 50\,\% bias improvement) plus the WP3.6
multi-port FIM (over-determined by 9$\times$) restore the
information rate the single-bin DFT cannot.

\subsection{Paramo-2023 detection $\to$ location extension}
The published Paramo et al. method DETECTS faults but does not
estimate location.  We extend it by sweeping $\alpha$ and selecting
the candidate that maximises the residual-covariance dominant
eigenvalue.  This is the only legitimate extension consistent with
the published statistic; the inverse direction (minimisation) does
not align with the detection-statistic intent.  The extension is
explicitly disclosed in the public module docstring and the
blind-review record.

\subsection{$\mu$-PMU two-ended via virtual-PMU}
The Cui-Weng-2020 method requires synchronised remote-end
measurements.  We provide them by emitting a virtual-PMU phasor
from the SAMBPS Digital Twin --- a parallel digital-twin instance
running at the remote substation that mirrors the system state
in real time.  This DT-as-virtual-PMU pattern is documented in the
WP4.5 brief; the competitor cannot tell that the source is a
digital twin.  In a deployment without a real $\mu$-PMU, the DT-
provided remote-end voltage is what a downstream consumer of the
Cui-Weng-2020 method would actually receive, so the comparison is
fair.

% =============================================================================
\section{Reproducibility statement}\label{sec:reproducibility}
% =============================================================================
All code, surrogate datasets, random seeds, and the exact
\texttt{outputs/phase4\_table3bis.csv} produced for this paper are
released under MIT licence in the public repository
\url{https://github.com/SAMBPS-DTaaS/HIF-TF-Locator}.  The full
benchmark is re-runnable from
\begin{verbatim}
  python run_faultloc_phase4_benchmark.py
\end{verbatim}
which produces \texttt{outputs/phase4\_table3bis.csv} and
\texttt{outputs/phase4\_figs/table3bis\_summary.pdf} directly.  The
public CNRS IEEE-34 HIF dataset~\cite{CNRS_2024_IEEE34} provides
an external real-world cross-check of the proposed method (CNRS
external validation reported in the companion conference
paper~\cite{Phase3conf2027}; full K05 measurement on the labelled
test slice is deferred behind a reviewer-licence gate).

% =============================================================================
\bibliographystyle{IEEEtran}
\bibliography{references}
% =============================================================================

\end{document}
